%!TEX root = ../main.tex

\chapter{Các Cấu Trúc Điều Khiển}
\hypertarget{localtoc\thechapter}{}
\localtoc

\begin{problem}
    Viết chương trình tìm số lớn nhất trong 3 số thực a, b, c
\end{problem}

\begin{problem}
    Viết chương trình nhập 2 số thực, kiểm tra xem chúng có cùng dấu hay không?
\end{problem}

\begin{problem}
    Viết chương trình giải và biện luận phương trình bậc nhất $ax + b = 0$
\end{problem}

\begin{problem}
    Nhập vào tháng của 1 năm. Cho biết tháng thuộc quý mấy trong năm
\end{problem}

\begin{problem}
    Tính $S(n) = 1^3 + 2^3 + … + n^3$
\end{problem}

\begin{problem}
    Tìm số nguyên dương n nhỏ nhất sao cho $1 + 2 + … + n > 10000$
\end{problem}

\begin{problem}
    Hãy sử dụng vòng lặp for để xuất tất cả các ký tự từ A đến Z
\end{problem}

\begin{problem}
    Viết chương trình tính tổng các giá trị lẻ nguyên dương nhỏ hơn N
\end{problem}

\begin{problem}
    Viết chương trình tìm số nguyên dương m lớn nhất sao cho 1 + 2 + … + m < N
\end{problem}

\begin{problem}
    In tất cả các số nguyên dương lẻ nhỏ hơn 100
\end{problem}

\begin{problem}
    Tìm ước số chung lớn nhất của 2 số nguyên dương
\end{problem}

\begin{problem}
     Tính tổng các số nguyên tố nhỏ hơn N (N nguyên dương)
\end{problem}

\begin{problem}
    Viết chương trình in ra tất cả các số lẻ nhỏ hơn 100 trừ các số 5, 7, 93
\end{problem}

\begin{problem}
    Viết chương trình nhập 3 số thực. Hãy thay tất cả các số âm bằng trị tuyệt đối của nó
\end{problem}

\begin{problem}
    Viết chương trình nhập giá trị x sau tính giá trị của hàm số
$ f(x) = 2x^2 + 5x + 9 khi x >= 5, f(x) = -2x^2 + 4x – 9 khi x < 5$
\end{problem}

\begin{problem}
    Viết chương trình nhập 3 cạnh của 1 tam giác, cho biết đó là tam giác gì
\end{problem}
\begin{problem}
    Lập chương trình giải hệ: $ax + by = c$, $dx + ey = f$. Các hệ số nhập từ bàn phím
\end{problem}
\begin{problem}
    Viết chương trình nhập vào 3 số thực. Hãy in 3 số ấy ra màn hình theo thứ tự tăng dần mà chỉ dùng tối đa 1 biến phụ
\end{problem}
\begin{problem}
    Viết chương trình giải phương trình bậc 2
\end{problem}
\begin{problem}
     Viết chương trình nhập tháng, năm. Hãy cho biết tháng đó có bao nhiêu ngày
\end{problem}
\begin{problem}
    Viết chương trình nhập vào 1 ngày (ngày, tháng, năm). Tìm ngày kế ngày vừa nhập (ngày, tháng, năm)
\end{problem}
\begin{problem}
    Viết chương trình nhập vào 1 ngày ( ngày, tháng, năm). Tìm ngày trước ngày vừa nhập (ngày, tháng, năm)
\end{problem}
\begin{problem}
    Viết chương trình nhập ngày, tháng, năm. Tính xem ngày đó là ngày thứ bao nhiêu trong năm
\end{problem}
\begin{problem}
    Viết chương trình nhập 1 số nguyên có 2 chữ số. Hãy in ra cách đọc của số nguyên này
\end{problem}
\begin{problem}
    Viết chương trình nhập 1 số nguyên có 3 chữ số. Hãy in ra cách đọc của số nguyên này
\end{problem}
\begin{problem}
    Viết hàm tính $S = \sqrt[n]{x}$
\end{problem}

\begin{problem}
    Viết hàm tính $S = x^y$
\end{problem}
\begin{problem}
    Viết chương trình in bảng cửu chương ra màn hình
\end{problem}
\begin{problem}
    Cần có tổng 200000 đồng từ 3 loại giấy bạc 1000 đồng, 2000 đồng, 5000 đồng. Lập chương trình để tìm ra tất cả các phương án có thể
\end{problem}
\begin{problem}
Viết chương trình in ra tam giác cân có độ cao h.
a. Tam giác cân đặc nằm giữa màn hình.
Ví dụ với h = 4\\
            *\\
          * * *\\
        * * * * *\\
      * * * * * * *\\
b. Tam giác cân rỗng nằm giữa màn hình.
Ví dụ với h = 4\\
*\\
* *\\
* *\\
* * * * * * *\\
c. Tam giác vuông cân đặc
Ví dụ với h = 4\\
*\\
* *\\
* * *\\
* * * *\\
d. Tam giác vuông cân rỗng
Ví dụ với h = 5\\
*\\
* *\\
* *\\
* *\\
* * * * *\\

\end{problem}
\begin{problem}
    Viết chương trình in ra hình chữ nhật có kích thước m x n\\
    \begin{enumerate}
        \item Hình chữ nhật đặc
        \item Hình chữ nhật rỗng
    \end{enumerate}
\end{problem}
\begin{problem}
    Lập chương trình tính sin(x) với độ chính xác 0.00001 theo công thức
 $Sin(x) = x – \frac{x^3}{3!} + \frac{x^5}{5!} + \cdots + (-1)^n \frac{x^{2n + 1}}{(2n + 1)!}$
\end{problem}
\begin{problem}
    Lam lai tat ca cac bai truoc bang HAM
\end{problem}


